\paragraph{Natural-shift breadth and interface utility.}
Across three separately locked natural-shift confirmations, MOSAIC releases 75/120 jobs (62.5\%) with 0/75 held-out violations: 20/20 BiasBios-Clinical, 15/60 ACS geographic, and 40/40 CINIC-10 image-origin jobs. The minimum leave-one-domain-out pooled release rate is 43.8\%. Relative to each selected four-token task interface, the median channel cost is 0.60 points on ACS and 0.60 points on CINIC-10; the maximum costs are 1.42 and 2.86 points. The supplement gives every seed and the full comparison boundary.

\paragraph{Prospective failure and estimated-source release.}
Two direct ACS interfaces that deployed on 2018 evidence, but that MOSAIC
rejected, were frozen for a 2023 Florida test after discovery in 2021 and
replication in 2022. Both exceed the $.40$ utility threshold empirically; one
is familywise confirmed, with error $.425$ and lower bound $.405$. A separate
proxy-source confirmation imputes source with $.841$ balanced accuracy and
certifies a nonconstant release from proxy-labeled target data. Its source
bounds are $.218/.279$, certified worst error is $.369$, and untouched
diagnostics are $.059$ and $.277$. Nested calibration folds abstain at 19,566
and 39,133 rows, then release at 58,700 and 78,267 rows.

\paragraph{What drives the bound.}
An exact decomposition over all 35 primary BiasBios and ACS releases finds
that residual shift, rather than multinomial sampling, is the larger source
and utility charge in every job. Median residual increments are $.191$ and
$.138$, while median sampling increments are below $5\times10^{-8}$. A
matched binary randomized-response baseline sets
$\epsilon_{\rm LDP}=.731$ so its contraction equals the $.35$ source
contract. It releases 0/35 interfaces at the $.40$ utility threshold; MOSAIC
releases 35/35 and lowers certified error in every job by a median $.155$.
